Figure~\ref{fig:models} and Table~\ref{tab:results} show the full model comparison. Over the tested lengths, the largest absolute discrepancy from the corpus is 12.91, 6.04, 2.52, and 1.56 percentage points for $Q_0,Q_1,Q_2,Q_3$, respectively. Thus three-letter context improves the central transition in this experiment, while the improvement from two to three letters of context is smaller than the improvement from independent letter frequencies to contextual models. This is an empirical comparison of these particular tables, not a monotonicity theorem about context length.

The corpus and all four English-model estimates first exceed $99\%$ at $N=18$ on the tested integer grid; the uniform-letter model first does so at $N=21$. At $N=18$, the corpus contains $M_{18}=2,772,123$ distinct passages and has 99.2075\% unique recovery, compared with 99.6368\% for the trigram model and 99.6028\% for the three-letter-context model. Nevertheless, corpus ambiguity is 0.7925\%, versus 0.3972\% under $Q_3$, approximately a factor of two. Agreement at the $99\%$ crossing therefore conceals a substantial relative discrepancy in the remaining ambiguity.

The largest estimated Monte Carlo standard error is 0.0165 percentage points, much smaller than the visible model--corpus discrepancies. Some long-length replicate estimates equal one because no ambiguous entries were observed; zero empirical variance there is not evidence of zero true ambiguity. Those zero ambiguity estimates are omitted on the logarithmic axis.

Figure~\ref{fig:interpretation} separates three goals. At $N=18$, optimal guessing succeeds with probability 99.5916\%, which exceeds unique identification. Remaining plaintext entropy is only 0.00838 bits, yet the mean key uncertainty even after revealing the plaintext is 38.55 bits. The repertoire has $\log_2 M_{18}=21.40$ bits of choice, or 1.189 bits per letter. The nominal $1.5$-bit asymptotic rate is not the measured rate of this finite list.
